\chapter{Why Rust for Data Science}
\label{ch:why-rust}

Most data scientists reach for Python. It is the default. The ecosystem is vast,
the learning curve is gentle, and nearly every ML paper ships a Python reference
implementation. So why would anyone choose Rust?

This chapter answers that question. We will set up a Rust environment, load a
real dataset, compute summary statistics, and produce a scatter plot --- all
before the end of the chapter. By the time you turn the page, you will have a
working project and a clear sense of what Rust brings to the table.


% ══════════════════════════════════════════════════════════════════════════════
\section{The Case for Rust in Data Science}
\label{sec:case-for-rust}

Rust is a systems language. It compiles to native code, manages memory without a
garbage collector, and enforces thread safety at compile time. These properties
matter for data science more than you might expect.

\subsection{Performance}

Rust programs run fast. There is no interpreter overhead, no garbage collector
pausing in the middle of a training loop, and no boxing of numeric types unless
you ask for it. The compiler applies aggressive optimizations: inlining,
autovectorization, loop unrolling, dead code elimination. Zero-cost abstractions
mean that iterators compile down to the same machine code as hand-written loops.

When your dataset grows from a CSV on your laptop to a hundred million rows on a
server, these properties stop being academic. A pipeline that takes four hours in
Python may finish in twenty minutes in Rust --- without reaching for C extensions
or Cython.

SIMD intrinsics are available when you need them, but in practice the
autovectorizer handles most numeric loops. The \texttt{ndarray} crate provides
$n$-dimensional arrays with BLAS-backed linear algebra, comparable to NumPy but
compiled ahead of time.

\subsection{Safety}

Rust's type system catches entire categories of bugs at compile time. There are
no null pointers --- \texttt{Option<T>} forces you to handle the absent case. There
are no data races --- the borrow checker prevents shared mutable state across
threads. There are no use-after-free bugs, no buffer overflows, no silent integer
truncation (unless you explicitly call \texttt{as}).

For data science, this means fewer midnight debugging sessions. If your pipeline
compiles, it will not segfault on row 50 million. Type errors in feature
engineering surface immediately, not after an hour of training.

\subsection{Reproducibility}

Cargo, Rust's build system, pins every dependency version in \texttt{Cargo.lock}.
A project that builds today will build identically six months from now, on any
platform, with the same compiler version. There is no \texttt{pip freeze} drift,
no conda environment corruption, no ``works on my machine.''

Cross-compilation is straightforward. A model trained on Linux can be deployed as
a single static binary on a minimal container image --- no Python runtime, no
shared libraries, no virtualenv.

\subsection{When Python is Still Better}

Rust is not always the right tool. Python excels at:

\begin{itemize}
  \item \textbf{Rapid prototyping.} A Jupyter notebook lets you explore data
    interactively. Rust's compile step, even if fast, adds friction to
    exploratory work.
  \item \textbf{Ecosystem breadth.} Python has libraries for everything:
    geospatial analysis, web scraping, bioinformatics, symbolic math. Rust's
    ecosystem is younger and narrower.
  \item \textbf{Team familiarity.} If your team knows Python and ships products
    in Python, switching languages has a real cost.
\end{itemize}

The pragmatic approach: prototype in Python, then rewrite the hot path in Rust.
Or, if you are starting fresh and performance matters from day one, start in
Rust. This book assumes the latter.


% ══════════════════════════════════════════════════════════════════════════════
\section{Setting Up Your Environment}
\label{sec:setup}

\subsection{Install Rust}

Install Rust through \texttt{rustup}, the official toolchain manager. On Linux
or macOS:

\begin{lstlisting}[language=bash, numbers=none]
curl --proto '=https' --tlsv1.2 -sSf https://sh.rustup.rs | sh
\end{lstlisting}

On Windows, download and run \texttt{rustup-init.exe} from
\url{https://rustup.rs}. Accept the defaults. After installation, open a new
terminal and verify:

\begin{lstlisting}[language=bash, numbers=none]
rustc --version
cargo --version
\end{lstlisting}

You should see version 1.75 or later. Any recent stable release will work.

\subsection{Create a New Project}

Cargo creates a new project with a single command:

\begin{lstlisting}[language=bash, numbers=none]
cargo new hello_ds
cd hello_ds
\end{lstlisting}

This generates the following structure:

\begin{lstlisting}[language=bash, numbers=none]
hello_ds/
  Cargo.toml
  src/
    main.rs
\end{lstlisting}

\texttt{Cargo.toml} is the project manifest. \texttt{src/main.rs} contains a
\texttt{main} function that prints ``Hello, world!'' --- we will replace it
shortly.

\subsection{Add Dependencies}

Open \texttt{Cargo.toml} and add two dependencies:

\begin{lstlisting}[language=toml]
[dependencies]
scry-learn = { version = "0.7", features = ["csv"] }
esoc-chart = "0.1"
\end{lstlisting}

\texttt{scry-learn} is the machine learning library we will use throughout this
book. The \texttt{csv} feature enables loading datasets from CSV files.
\texttt{esoc-chart} provides visualization through a grammar-of-graphics API.

\subsection{Essential Cargo Commands}

Four commands cover most of your workflow:

\begin{lstlisting}[language=bash, numbers=none]
cargo build          # compile the project
cargo run            # compile and run
cargo test           # run tests
cargo doc --open     # generate and open API documentation
\end{lstlisting}

Cargo fetches dependencies automatically on the first build. Subsequent builds
are incremental and fast.


% ══════════════════════════════════════════════════════════════════════════════
\section{Hello, Data Science}
\label{sec:hello-ds}

Let us do something real. We will load the Iris dataset, print summary
statistics, and produce a scatter plot. This is the ``Hello, World'' of data
science, and it will exercise the core workflow: load data, inspect it,
visualize it.

\subsection{Loading the Dataset}

The Iris dataset contains 150 measurements of three species of iris flowers.
Each row has four features (sepal length, sepal width, petal length, petal
width) and a class label (setosa, versicolor, or virginica).

\texttt{scry-learn} provides a \texttt{Dataset} struct that holds feature
matrices and target vectors together. The \texttt{from\_csv} constructor reads a
CSV file and designates one column as the target:

\begin{lstlisting}[language=rust]
use scry_learn::dataset::Dataset;

let ds = Dataset::from_csv("datasets/iris/iris.csv", "species")
    .expect("failed to load iris.csv");
\end{lstlisting}

The second argument, \texttt{"species"}, names the target column. All remaining
columns become features. Categorical targets are automatically encoded as
integers, and the original class labels are stored in
\texttt{ds.class\_labels}.

\subsection{Summary Statistics}

Call \texttt{describe()} to print a summary table:

\begin{lstlisting}[language=rust]
ds.describe();
\end{lstlisting}

This prints a table like the following:

\begin{lstlisting}[numbers=none]
              count   mean    std     min     25%     50%     75%     max
sepal_length  150     5.843   0.828   4.300   5.100   5.800   6.400   7.900
sepal_width   150     3.054   0.434   2.000   2.800   3.000   3.300   4.400
petal_length  150     3.759   1.764   1.000   1.600   4.350   5.100   6.900
petal_width   150     1.199   0.763   0.100   0.300   1.300   1.800   2.500
\end{lstlisting}

You immediately see the scale of each feature. Sepal length ranges from 4.3 to
7.9; petal width from 0.1 to 2.5. This kind of inspection is the first step in
any analysis.

\subsection{Making a Scatter Plot}

The \texttt{esoc-chart} crate provides an express API for common plots. To make
a scatter plot of sepal length versus sepal width, colored by species, we need
the feature columns and the class labels:

\begin{lstlisting}[language=rust]
use esoc_chart::express::scatter;

// Extract feature columns
let sepal_length = ds.column("sepal_length");
let sepal_width = ds.column("sepal_width");

// Map encoded target back to species names
let species: Vec<String> = ds.target.iter()
    .map(|&t| {
        ds.class_labels.as_ref().unwrap()[t as usize].clone()
    })
    .collect();

// Build the scatter plot
let svg = scatter(&sepal_length, &sepal_width)
    .color_by(&species)
    .title("Iris: Sepal Length vs Width")
    .x_label("Sepal Length")
    .y_label("Sepal Width")
    .to_svg();

std::fs::write("iris_scatter.svg", &svg)
    .expect("failed to write SVG");

println!("Wrote iris_scatter.svg");
\end{lstlisting}

This produces an SVG file with three color-coded clusters. Setosa separates
cleanly from the other two species along sepal width. Versicolor and virginica
overlap --- a pattern we will revisit when we build classifiers in
Chapter~\ref{ch:classification}.

\subsection{The Complete Program}

Here is the full \texttt{src/main.rs}:

\begin{lstlisting}[language=rust, caption={src/main.rs --- Hello Data Science}]
use scry_learn::dataset::Dataset;
use esoc_chart::express::scatter;

fn main() {
    // Load the Iris dataset
    let ds = Dataset::from_csv(
        "datasets/iris/iris.csv", "species"
    ).expect("failed to load iris.csv");

    // Print summary statistics
    ds.describe();

    // Extract feature columns
    let sepal_length = ds.column("sepal_length");
    let sepal_width = ds.column("sepal_width");

    // Map encoded target back to species names
    let species: Vec<String> = ds.target.iter()
        .map(|&t| {
            ds.class_labels
                .as_ref()
                .unwrap()[t as usize]
                .clone()
        })
        .collect();

    // Build a scatter plot
    let svg = scatter(&sepal_length, &sepal_width)
        .color_by(&species)
        .title("Iris: Sepal Length vs Width")
        .x_label("Sepal Length")
        .y_label("Sepal Width")
        .to_svg();

    std::fs::write("iris_scatter.svg", &svg)
        .expect("failed to write SVG");

    println!("Wrote iris_scatter.svg");
}
\end{lstlisting}

Run it:

\begin{lstlisting}[language=bash, numbers=none]
cargo run
\end{lstlisting}

Cargo compiles the project, fetches dependencies, and executes the binary. You
should see the summary statistics table followed by the confirmation message.
Open \texttt{iris\_scatter.svg} in a browser to see the plot.


% ══════════════════════════════════════════════════════════════════════════════
\section{What Comes Next}
\label{sec:what-next}

This book is organized in four parts.

\textbf{Part~I: Foundations} (Chapters~\ref{ch:why-rust}--\ref{ch:preprocessing})
covers the prerequisites: working with tabular data, visualization, and data
cleaning. These chapters establish the patterns you will use in every subsequent
chapter.

\textbf{Part~II: Supervised Learning}
(Chapters~\ref{ch:linear-regression}--\ref{ch:gradient-boosting}) builds models
that predict outcomes from labeled data. We start with linear regression, move
through decision trees and random forests, and finish with gradient boosting.
Each chapter includes model evaluation and practical advice on when to use which
algorithm.

\textbf{Part~III: Unsupervised Learning and Beyond}
(Chapters~\ref{ch:clustering}--\ref{ch:svm-knn}) covers clustering, anomaly
detection, text analysis, neural networks, and support vector machines. These
methods find structure in data without explicit labels, or operate in domains
where the supervised framing is less natural.

\textbf{Part~IV: Production and Practice}
(Chapters~\ref{ch:hyperparameter}--\ref{ch:case-study}) addresses the concerns
that arise when models leave the notebook: hyperparameter tuning,
explainability, calibration, ensembling, and a complete end-to-end case study
that ties everything together.

Turn the page. We have data to work with.
